\section{Evaluation}\label{sec:eval}

We establish the practicality of the scheme by integrating the protocol into
a mobile wallet and benchmarking its performance on both iOS and Android.
Figures~\ref{fig:msgSeqIssue} and~\ref{fig:msgSeqTransfer} show our
implementation of the withdrawal and payment protocols of
Sections~\ref{sec:withdrawal} and~\ref{sec:payments}. Four entities are
involved: a mobile phone, a secure element, an intermediary, and a central
bank. We focus our evaluation on the mobile phone and the secure element, to
establish whether they can implement the protocol efficiently given their
hardware and storage limitations; the intermediary and the central bank need
no such evaluation, since they are expected to run on servers or in the
cloud. The mobile phone and secure element take on three different personas
across the protocol: the holder, who receives funds during withdrawal, and
the payer and payee, who exchange possession of a token during payment.

\subsection{Implementation}

\paragraph{Secure element.} A fundamental aspect of a realistic benchmark is
the choice of secure element. We consider four options: software emulation
of the secure element in the wallet app; a TEE-based secure element in the
mobile platform; an HSM-based secure element in the mobile platform; and an
external secure element. We discard software emulation, since it would
present unrealistically good performance by ignoring both the
resource-constrained environment of a real secure element and the
communication overhead of talking to one. We discard TEE-based secure
elements as insufficiently secure, given their wider attack surface and
history of compromise~\cite{TEEsBad}. We discard HSM-based secure elements
because those integrated into Apple and Android platforms do not implement
our scheme and are not extensible. We therefore adopt an external secure
element, implemented on \emph{smart cards}.

A smart card is a physical card embedded with an integrated circuit that can
process and securely store data, including a dedicated security component
that protects sensitive data and cryptographic operations. We use
Javacard~\cite{javacard}, a Java-based framework for smart cards, in which
the protocol running on the secure element is implemented as an
\emph{applet}, a small application supported by a subset of the Java
runtime. Javacard applets communicate through Application Protocol Data
Units (APDUs), which identify the applet function that should process a
given request; the Javacard framework supports elliptic curve operations.

Our smart card implementation exposes two functions, \emph{receive} and
\emph{sign}. \emph{Receive} implements the secure element's part of the
protocols in Sections~\ref{sec:stepW1W2} and~\ref{sec:stepS1}, which are the
same protocol: sampling randomness and generating a Pedersen commitment.
\emph{Sign} implements the smart card's part of the protocol in
Section~\ref{sec:stepS2S3S4}, described in more detail in
Appendix~\ref{sec:joint-proof}.

Recall that, with the smart card's help, the payer produces a signature of
knowledge over two commitments, the payer's and the payee's ($C_i$ and
$C_{i+1}$ in Section~\ref{sec:stepS2S3S4}). To prevent double spending, the
smart card is programmed to sign only $(v, C_i, C_{i+1})$ triples such that
$C_i$ is the commitment the card generated when it received the token, and
to sign against only a single destination commitment $C_{i+1}$ for that
token once it is spent. The first constraint ensures the card signs off
payments only for the token last received, avoiding re-spending of past
ones; the second ensures the current token can be transferred to only one
recipient. To tolerate errors on the NFC channel between the smart card and
the mobile phone, which are not uncommon, we permit multiple invocations of
\emph{sign} provided the supplied destination commitment $C_{i+1}$ never
changes, all of which generate valid signatures of knowledge for the same
payer/payee commitment pair; the card is reset by invoking \emph{receive},
which deletes the previous commitment and generates a new one. We introduce
the concept of a \emph{slot}, each able to hold one pair of payer/payee
commitments, to let a single card handle multiple tokens; the protocols
above simply need to be augmented with a slot identifier.

\paragraph{Mobile wallet.} The mobile wallet is an app that implements the
functionality needed for a user to withdraw and transfer tokens. It handles
communication with the intermediary, which runs as a cloud service, between
users during a payment, and with the smart card. System parameters and
credentials for each participant are generated by a registration authority,
also exposed as a cloud service.

To implement withdrawal and payment, the wallet exposes \emph{reqGen} (the
holder device's part of the protocol in Section~\ref{sec:stepW1W2}),
\emph{unblindAndVerify} (the protocol in Section~\ref{sec:stepW4}), and
\emph{makeTokens} (the protocol in Section~\ref{sec:stepW5}) for withdrawal;
and \emph{prepare} (the encryption and correctness proof of
Section~\ref{sec:stepS2S3S4}), \emph{finalize} (the payer device's part of
the protocol in Section~\ref{sec:stepS2S3S4}), and \emph{verify} (the
protocol in Section~\ref{sec:stepS5}) for payment.

\begin{figure*}[tp]
  \begin{minipage}{.48\textwidth}
    \centering
    \includegraphics[trim={1.5cm 10.9cm 1.8cm 1cm},clip,width=0.8\columnwidth]{measurements/issue}
    \caption{Sequence diagram for token withdrawal.\label{fig:msgSeqIssue}}
  \end{minipage}
  \quad \quad
  \begin{minipage}{.48\textwidth}
    \centering
    \includegraphics[trim={1.9cm 10.9cm 1.5cm 1cm},clip,width=0.8\columnwidth]{measurements/spend}
    \caption{Sequence diagram for token payment.\label{fig:msgSeqTransfer}}
  \end{minipage}
\end{figure*}

\subsection{Protocols}

The message sequence flow for withdrawal and payment is shown in
Figures~\ref{fig:msgSeqIssue} and~\ref{fig:msgSeqTransfer}.

\paragraph{Withdraw.} Withdrawing a token involves local invocations of the
mobile function, NFC communication between the mobile phone and the smart
card, and a single remote call to the intermediary, which in turn contacts
the central bank (omitted from the message sequence diagram for space). The
protocol starts with the holder scanning her card to obtain values $s$ and
$c$, using these to generate a withdrawal request (steps \textbf{W1} and
\textbf{W2}, Section~\ref{sec:stepW1W2}), and issuing it to the REST
endpoint of her intermediary over HTTPS. The intermediary interacts directly
with the central bank, which runs the protocol of Section~\ref{sec:stepW3}
and returns a response to the holder (step \textbf{W3}). The holder then
runs the protocols of Sections~\ref{sec:stepW4} and~\ref{sec:stepW5} to
create a token, which is stored for future use (steps \textbf{W4} and
\textbf{W5}).

\paragraph{Payment.} Before the offline payment protocol of
Figure~\ref{fig:msgSeqTransfer} can start, payer and payee negotiate the
terms of the payment and set up a BLE channel. The payee displays a request,
base64-encoded as a QR code, containing the payment terms and a BLE service
identifier, and starts a BLE GATT server advertising that service; under the
service, a \emph{Notify} characteristic sends data to the payer, and a
\emph{Write without response} characteristic receives data from the payer.
The QR code also carries the cryptographic material for a secure BLE
exchange based on Diffie-Hellman. The payer scans the QR code to retrieve
the BLE information and payment details and, if she agrees, completes the
Diffie-Hellman exchange over BLE and confirms the payment terms; all
subsequent messages, including this one, are encrypted and CBOR-encoded.

The protocol proper begins with the payee scanning her card (step
\textbf{S1}) to obtain values $s$ and $C$ (Section~\ref{sec:stepS1}) and
sending $C$ to the payer over BLE. The payer prepares the payment (step
\textbf{S2}) by generating the audit information, the encryption of her
identifier together with its proof of correctness
(Section~\ref{sec:stepS2S3S4}), and then her device and secure element
engage in the joint proof of knowledge protocol of
Appendix~\ref{sec:joint-proof} (steps \textbf{S3} and \textbf{S4}). The
resulting payment $\py$ is sent to the payee over BLE, who verifies the
received token (step \textbf{S5}) using the verification algorithm of
Section~\ref{sec:stepS5} and stores it upon success.

\subsection{Results}

We implement the protocol on Android (Samsung Galaxy A34, Android 14) and
Apple (iPhone 13 mini, iOS 17.6.1) platforms, both running a crypto library
built with Go 1.21.8. We benchmark two versions of the protocol, with and
without auditing, i.e. with and without verifiable encryption of the
payer's identifier for the auditor, to clearly show the overhead auditing
imposes. The smart card is an NXP 1ID white plastic card with a Smart MX
D600 chip (400\,KiB of available memory) running JCOP 4.5 OS with NXP's
JCOPx extensions at version 1.1.4; card and phone communicate over NFC.

Figures~\ref{fig:perfWithdraw} and~\ref{fig:perfPayment} show the
performance of the no-audit version of withdrawal and payment. Withdrawal
and the payer side of a payment complete in under $0.5\,\mathrm{s}$ on both
platforms; the payee's running time on payment instead depends on the
length of the payment chain, since the payee must verify its correctness,
and stays below 1 second for chains of up to 32 hops. In both protocols,
the dominant steps are \textbf{W1} for withdrawal and \textbf{S1} and
\textbf{S3} for payment, dominated by the point addition and scalar
multiplication needed to compute the Pedersen commitment on the card; step
\textbf{S5} is negligible for short chains.

Figure~\ref{fig:audit} shows the overhead auditing imposes on the payment
protocol for both payer and payee. Auditing has no impact on the payer, and
imposes an overhead of between 10 and 20 percent on the payee's running
time; we speculate that the lower overhead observed for shorter chains
reflects more moderate memory usage when verifying short chains, masking
the effects of Go's garbage collector for longer ones.

\begin{figure*}[tp]
  \begin{minipage}{.3\textwidth}
    \centering
    \includegraphics[trim={1cm 0.8cm 1cm 1cm},clip,width=\textwidth]{measurements/withdraw.png}
    \caption{Performance of withdrawal, broken into component
    steps.\label{fig:perfWithdraw}}
  \end{minipage}
  \quad
  \begin{minipage}{.3\textwidth}
    \centering
    \includegraphics[trim={1cm 0.8cm 1cm 1cm},clip,width=\textwidth]{measurements/transfer.png}
    \caption{Performance of payment, broken into
    steps.\label{fig:perfPayment}}
  \end{minipage}
  \quad
  \begin{minipage}{.3\textwidth}
    \centering
    \includegraphics[trim={1cm 0.8cm 1cm 1cm},clip,width=\textwidth]{measurements/overhead.png}
    \caption{Overhead of auditing on the payment protocol.\label{fig:audit}}
  \end{minipage}
\end{figure*}

\paragraph{Message size.} With reference to Figures~\ref{fig:msgSeqIssue}
and~\ref{fig:msgSeqTransfer}, and using the configuration above, $c$ is 171
bytes, $req$ is 459 bytes, $resp$ is 72 bytes, and $\mathfrak{p}$ is 1382
bytes, growing by roughly 1\,KiB for each additional transfer. This linear
growth is not a concern in practice: offline payment chains are expected to
be short, and are often explicitly limited, for instance under Annex 1 of
the digital euro documentation~\cite{ECBAnnex1}.

\paragraph{State storage requirements.} The issuer must store two lists: all
issued tokens and all deposited tokens. The list of deposited tokens is
needed to detect double spending and can never be pruned outright, though
mechanisms such as logical or wall-clock timestamps that bound a token's
validity span allow it to be pruned once that span elapses. The list of
issued tokens is kept only to support monetary policy, and is not required
by the protocol itself. Intermediaries typically store user accounts to
determine reserve requirements, again not a protocol requirement, and
users store the payment chains of any unspent tokens they hold.
